\documentclass[10pt]{article}
\usepackage[a4paper,margin=25mm]{geometry}
\usepackage{amsmath,amssymb}
\usepackage[T1]{fontenc}
\usepackage{hyperref}
\title{A cubic-box counterexample to Conjecture 425}
\author{Local AI-assisted solution draft}
\date{3 October 2026}
% Compact consistent title for a short mathematical note.
\makeatletter
\renewcommand{\maketitle}{\begin{center}
{\Large\bfseries\@title\par}\vspace{0.6em}
{\small\@author\par\@date}
\end{center}\vspace{0.5em}}
\makeatother
\setlength{\emergencystretch}{3em}
\usepackage{url}
\begin{document}
\maketitle
\paragraph{Source and interpretation.}
Conjecture \texttt{00000000425} in The Last Math Competition, at commit
\path{efab34b80a63963991d6c7ed625442a89a328a44}, asserts that the
coefficient sequence in the diagonal specialization of the boxed plane
partition generating function is always log-concave. We use the standard
volume generating polynomial, with the diagonal specialization meaning equal
box side lengths. The counterexample is the $2\times2\times2$ box.

\paragraph{Definitions.}
A plane partition in this box is an array
$\left(\begin{smallmatrix}a&b\\c&d\end{smallmatrix}\right)$ with entries in
$\{0,1,2\}$, weakly decreasing along both rows and columns:
$a\ge b$, $a\ge c$, $b\ge d$, and $c\ge d$.
Its volume is $a+b+c+d$. Let $c_k$ count such arrays of volume $k$.
Log-concavity requires $c_k^2\ge c_{k-1}c_{k+1}$ at every interior index.

\paragraph{Counterexample and proof.}
There is exactly one array of volume zero, namely the zero array, and exactly
one of volume one, with a single $1$ in the top-left entry. The three arrays
of volume two are
\[
\begin{pmatrix}2&0\\0&0\end{pmatrix},\qquad
\begin{pmatrix}1&1\\0&0\end{pmatrix},\qquad
\begin{pmatrix}1&0\\1&0\end{pmatrix}.
\]
These are exhaustive: either one entry equals $2$, which by monotonicity
must be the top-left entry, or two entries equal $1$, giving exactly the
two displayed order ideals. Consequently
\[
c_0=1,\qquad c_1=1,\qquad c_2=3,\qquad
c_1^2=1<3=c_0c_2.
\]
Thus the asserted universal log-concavity is false. It is unnecessary to
consider the additional equality characterization after this first clause
has failed.

\paragraph{Independent polynomial check.}
The complete polynomial, obtained both by direct enumeration and by the
MacMahon product, is
\[
\prod_{i,j,k=1}^{2}\frac{1-q^{i+j+k-1}}{1-q^{i+j+k-2}}
=1+q+3q^2+3q^3+4q^4+3q^5+3q^6+q^7+q^8.
\]
The accompanying Python verifier independently computes the product as an
exact integer formal power series and enumerates the arrays.

\paragraph{Formal verification and scope.}
The self-contained Lean 4.19.0 file \texttt{Main.lean} defines all four entries
as \texttt{Fin 3}, proves the 81-array enumeration is complete and has no
duplicates, proves the filtered membership predicate is exactly the
plane-partition condition, and computes every coefficient. The theorem
\path{Conjecture425.conjecture425_counterexample} proves the negation of
log-concavity for this genuine diagonal box. All calculations use kernel
reduction; there are no proof placeholders, custom axioms, or native decision
procedures. The formal proof follows the direct counting definition; the
MacMahon identity is an independent check, not an assumed Lean premise.

\paragraph{Alternative reading of ``coefficientwise.''}
If the intended claim instead compares consecutive polynomials $F_n(q)$
for $n\times n\times n$ boxes, it is false as well. The completely filled
box is the unique term of maximum volume, so $F_n$ is monic of degree
$n^3$. At $n=2$, the polynomial $F_2(q)^2-F_1(q)F_3(q)$ has coefficient
$-1$ at $q^{28}$, because the two product degrees are $16$ and $28$.
This additional argument covers that reading in ordinary mathematics;
the Lean artifact formalizes the scalar coefficient-sequence reading above.

\end{document}
